%% CV_DPN.tex — Dipayan Sarkar
\documentclass[11pt, a4paper]{moderncv}

% \moderncvstyle{classic}
\moderncvstyle{banking}   % or casual, fancy, oldstyle
% \moderncvstyle{casual}   % or casual, fancy, oldstyle
% \moderncvstyle{fancy}   % or casual, fancy, oldstyle
% \moderncvstyle{oldstyle}   % or casual, fancy, oldstyle


\moderncvcolor{black}
% \moderncvcolor{orange}    % or green, red, purple, grey, black

% Force first name and last name to the same color (keep last name bold)
\renewcommand*{\firstnamestyle}[1]{{\fontsize{34}{36}\bfseries\upshape\color{black}#1}}
\renewcommand*{\lastnamestyle}[1]{{\fontsize{34}{36}\bfseries\upshape\color{black}#1}}

% Add vertical space between the name and the address/contact links in the header
\makeatletter
\patchcmd{\makehead}{\\\addressfont\color{addresscolor}}{\\[1em]\addressfont\color{addresscolor}}{}{%
  \PackageWarning{CV_DPN}{Could not patch \string\makehead\space for name spacing}}
\makeatother

% \renewcommand{\namefont}{\Huge\bfseries\color{black}}
% \renewcommand{\titlefont}{\large\mdseries\upshape\color{black}}

\usepackage[scale=0.85, top=1.8cm, bottom=1.5cm]{geometry}
\usepackage{multicol}
\usepackage{lastpage}   % reliable \pageref{LastPage}; moderncv's own lastpage label never fires here
% \usepackage{mathptmx}   # times new roman font
% \usepackage{palatino}                                         
% \usepackage{fontspec}   
\usepackage{lmodern}
% fontawesome icons are provided by moderncv (fontawesome6); do not load fontawesome5 (clashes)
\newcommand{\huggingfaceicon}{\raisebox{-0.2em}{\includegraphics[height=0.99em]{hf_logo.png}}}


% Personal inf
\name{Dipayan}{Sarkar}
\address{Department of Bioinformatics, University of North Bengal}{Siliguri - 734013, West Bengal, India.}{}
\phone[mobile]{+91 9735490241}
\email{dipayanskr@gmail.com}
\homepage{dipayansarkar.com}
\social[github][github.com/dipayanskr]{github.com/dipayanskr}
\social[googlescholar][scholar.google.com/citations?user=K-zTRIEAAAAJ&hl=en]{Google Scholar}
% \social[linkedin][www.linkedin.com/in/dipayan-sarkar-b5a281179]{dipayansarkar}
\social[orcid][orcid.org/0000-0002-0765-5689]{0000-0002-0765-5689}
\extrainfo{%
  \emailsymbol\emaillink{rs\_dipayans@nbu.ac.in}%
  \makeheaddetailssymbol
  \huggingfaceicon~\href{https://huggingface.co/dipayanskr}{huggingface.co/dipayanskr}%
}










% ======================================================================
\begin{document}

\makecvtitle

% Running footer: name (left) + page number (right) on continuation pages only.
% moderncv puts its page number inside \fancypagestyle{plain}, so both fields must be
% redefined together here -- a bare \fancyfoot[l] would drop the page number.
\fancypagestyle{plain}{%
  \fancyhf{}%
  \fancyfoot[l]{\ifnum\value{page}>1\relax
    \color{color2}\pagenumberfont\strut Dipayan Sarkar\fi}%
  \fancyfoot[r]{\color{color2}\pagenumberfont\strut\thepage/\protect\NoHyper\pageref{LastPage}\protect\endNoHyper}%
}
\pagestyle{plain}

% ---------- Research Summary ----------
\section{Research Summary}
\cvitem{}{%
	PhD researcher in bioinformatics and computational biology developing deep learning methods for protein sequence modeling. Work spans sequence-based protein-protein interaction prediction using attention-based architectures with transfer learning from pretrained protein language models, generative modeling of protein binders through autoregressive and diffusion-based approaches. Also builds web platforms that make trained models usable for inference and interaction-network visualization.
	% Ongoing exploration includes LLM fine-tuning, reinforcement learning for sequence optimization, and discrete diffusion for biological sequence generation.
}

% ---------- Work Experience ----------

% ---------- Education ----------
\section{Education}

\cventry{Expected Submission 2027}{Ph.D. in Bioinformatics}{University of North Bengal, Department of Bioinformatics}{India}{}{%
	Supervisor: Dr.\ Chiranjib Sarkar, Assistant Professor, Department of Bioinformatics\\
	Focus: Framework Development for Predicting Protein-Protein Interaction Networks Using Deep Learning Algorithms on Protein Sequences.
}

\cventry{2021}{M.Sc. in Botany}{University of North Bengal, Department of Botany}{India}{77.25\% - First Class}{%
	Specialization: Biochemistry\\
	Dissertation: ``A review on the responses of aquatic macrophytes to phenolic compounds with special reference to their tolerance strategies and future implications.''\\
	Supervisor: Dr.\ Swarnendu Roy, DRS Department of Botany.
}

\cventry{2019}{B.Sc. (Hons.) in Botany}{Ananda Chandra College, University of North Bengal}{India}{63.50\% - First Class}{}

% \cventry{2016}{Higher Secondary (Class XII)}{Haldibari High School (H.S.)}{West Bengal}{80.00\% -- First Class}{%
%   WBCHSE | Subjects: Physics, Chemistry, Mathematics, Biology, English, Bengali
% }

% \cventry{2014}{Secondary (Class X)}{Haldibari High School (H.S.)}{West Bengal}{91.58\% -- First Class}{%
%   WBBSE | Subjects: Mathematics, Physical Science, Life Science, History, Geography, English, Bengali
% }

\section{Research Experience}

\cventry{2022--Present}{PhD Research Scholar}{Department of Bioinformatics, University of North Bengal}{Siliguri, India}{}{%
	Department of Bioinformatics
	\begin{itemize}
		\item[\textbullet] Deep learning architectures for sequence-based protein-protein interaction prediction using Attention-based hybrid architecture with transfer learning.
		\item[\textbullet] Sequence based generative modeling of protein binder generation using autoregressive and diffusion-based approaches.
		\item[\textbullet] Development of web-based platforms for model inference and protein interaction network visualization
	\end{itemize}
}




% ---------- Publications ----------
\section{Publications}

\cvitem{2025}{%
	\textbf{Sarkar, D.}, \& Sarkar, C. \textbf{AttnSeq-PPI}: Enhancing protein-protein interaction network prediction using transfer learning-driven hybrid attention. \textbf{\textit{Biochimica et Biophysica Acta (BBA)---Proteins and Proteomics}}, \textbf{1874}(1), 141102. \href{https://doi.org/10.1016/j.bbapap.2025.141102}{doi:10.1016/j.bbapap.2025.141102}
}
\cvitem{2026}{%
	\textbf{Sarkar, D.}, \& Sarkar, C. \textbf{ARACoFusion}: Uncertainty-aware calibrated deep learning for protein-protein interaction network prediction in \textit{Arabidopsis thaliana}. \textbf{bioRxiv} preprint. \href{https://www.biorxiv.org/content/10.64898/2026.05.22.727120v1}{doi:10.64898/2026.05.22.727120}
}

\cvitem{2026}{%
\textbf{Sarkar, D.}, \& Sarkar, C. \textbf{MoE-Bind}: Guiding De Novo Protein Binder Generation with Sparse Experts. \textbf{bioRxiv} preprint. \href{https://www.biorxiv.org/content/10.64898/2026.06.13.732043v1}{doi:10.64898/2026.06.13.732043}
}

\cvitem{2026}{%
	\textbf{Sarkar, D.}, Bardhan, K., \& Sarkar, C. \textbf{Deep-Interact Studio}: An Interactive Deep Learning Model Building Platform for Biomolecular Interaction Prediction. \textbf{bioRxiv} preprint. \href{https://www.biorxiv.org/content/10.64898/2026.07.02.736034v1}{doi:10.64898/2026.07.02.736034}
}





% ---------- Skills ----------
% \newpage

% ---------- Software & Tools ----------
\section{Research Software Developed}

\subsection{Web Servers}

\cvitem{Deep-Interact Studio}{%
	Interactive web platform for building, training, and comparing custom deep learning models for protein--protein, drug--target, RNA--protein, and protein--DNA interaction prediction, with integrated interpretability.
	\newline\href{https://deepinteract.compbiosysnbu.in/}{\texttt{deepinteract.compbiosysnbu.in}}
}

\cvitem{AttnSeq-PPI}{%
	Deep learning-based protein--protein interaction prediction platform.
	\newline\href{https://attnseqppi.compbiosysnbu.in/}{\texttt{attnseqppi.compbiosysnbu.in}}
}

% \cvitem{ARACoFusion-PPI}{%
%   Arabidopsis-specific PPI prediction and network analysis tool.
%   \newline\href{https://aracofusion.compbiosysnbu.in/}{\texttt{aracofusion.compbiosysnbu.in/}}
% }

\subsection{Packages}

\cvitem{EGRNi}{%
	R package for ensemble gene regulatory network inference, combining multiple inference
	methods to improve reconstruction accuracy of gene regulatory networks.
	\newline Sarkar, C., \textbf{Sarkar, D.}, Parsad, R., \& Mishra, D. \textit{CRAN}, v0.1.6
	\newline\href{https://cran.r-project.org/package=EGRNi}{\texttt{cran.r-project.org/package=EGRNi}}
}

% ---------- Teaching & Mentoring ----------
% \section{Teaching \& Mentoring}

% \cvitem{2022-2024}{%
% 	\textbf{M.Sc. Bioinformatics (CBCS)} - University of North Bengal. Independent theory and laboratory instruction in three papers (9 credits), as assigned by the course in-charge:
% 	\begin{itemize}
% 		\item[\textbullet] \textit{Basic Computer Applications} (Core, theory and practical, 4 cr; Sem.\ I, 2022-23 \& 2023-24)
% 		\item[\textbullet] \textit{Computer Programming for Bioinformatics} (Core, theory and practical, 4 cr; Sem.\ II, 2022-23)
% 		\item[\textbullet] \textit{Basic Mathematics and Statistics for Biology} (DSE practical, 1 cr; Sem.\ I, 2022-23 \& 2023-24)
% 	\end{itemize}
% }

% \cvitem{2025}{%
% 	\textbf{Dissertation mentoring} - mentoring of an M.Sc. Bioinformatics dissertation student (Central University of South Bihar) on \textit{Deep Learning-Based Model for Prediction of Targeted Sequencing Depth Using Probe Sequences}; submitted May 2025.
% }

\section{Skills}

\cvitem{Programming}{Python (PyTorch), R, C }
\cvitem{Deep Learning}{Transformers, attention mechanisms, transfer learning, pre-training and finetuning.}
\cvitem{Generative AI}{Autoregressive generative models for de novo protein and protein binder design.}
\cvitem{Scalable Training}{GPU-based training, parallel training on HPC environments}
\cvitem{Bioinformatics}{Biological sequence analysis, protein-protein interaction and biomolecular network prediction.}
\cvitem{Scientific Computing}{NumPy, Pandas, Matplotlib, Seaborn.}
\cvitem{Deployment}{Nginx, Docker, Flask, FastAPI, Linux server deployment, model-backed web applications.}
\cvitem{Tooling}{Git, GitHub Actions (CI), \texttt{uv}, Conda.}


% ---------- Talks & Presentations ----------
\section{Talks \& Presentations}

\cvitem{2025}{%
	\textbf{A Webtool for Transfer Learning Based Protein-Protein Interaction Network Prediction Using Hybrid Attention} - Oral presentation, \textit{Anusandhan 2025} (SARANSH, Wiley-sponsored), Bose Institute, Kolkata, India, 7 March 2025.
}


% ---------- Awards & Certificates ----------
\section{Awards}

% \cvitem{2021}{%
%   Qualified \textbf{CSIR-NET (JRF)} in Life Sciences - June 2021, \textbf{All India Rank: 216}
% }
\cvitem{2021}{%
	\textbf{CSIR-UGC NET with JRF}, Life Sciences - June 2021
	\newline\textit{All India Rank 216 \ $\vert$ \ Junior Research Fellowship, CSIR-HRDG, Government of India}
}
\cvitem{2022}{%
	Qualified \textbf{GATE} (Graduate Aptitude Test in Engineering) - Life Sciences (XL)
}

% \section{Fellowships}
% \cvitem{2021}{%
% 	\textbf{CSIR-UGC Junior Research Fellowship (JRF)} - Council of Scientific \& Industrial Research, Government of India. Awarded on qualifying CSIR-UGC NET (Life Sciences, June 2021).
% }

% ---------- Grants & Computing Resources ----------
\section{Grants \& Computing Resources}

% \cvitem{2025}{%
% 	\textbf{IndiaAI Compute Initiative} - Awarded GPU cloud compute (NVIDIA L4) under the IndiaAI Mission for the project \textit{Deep learning-based framework for protein-protein interaction network prediction} (Project ID: P1-S2025070964, Sep 2025 -- Aug 2026). Funded under CSIR-UGC SRF Fellowship.
% }
% \cvitem{2025}{%
% 	\textbf{Google TRC (TPU Research Cloud) Award} - Granted free access to Google Cloud TPUs (v4, v5e, v6e) for machine learning research.
% }
\cvitem{2025}{%
	\textbf{IndiaAI Compute Initiative} - Awarded GPU cloud compute (NVIDIA L4 GPU) under the IndiaAI Mission for the project \textit{Deep learning-based framework for protein-protein interaction network prediction} (Project ID: P1-S2025070964, Sep 2025 - Aug 2026).
}
\cvitem{2025}{%
	\textbf{Google TRC (TPU Research Cloud) Award} - Awarded access to GCP - TPUs (v4, v5e, v6e) for machine learning research.
}



% #############


% #############
% #############



\end{document}